\documentclass[12pt,a4paper]{article}
\usepackage[margin=25mm, headheight=15pt, headsep=10pt]{geometry}
\usepackage{fontspec}
\usepackage[table]{xcolor} % Note: table option is required for rowcolors
\usepackage{titlesec}
\usepackage{tocloft}
\usepackage{fancyhdr}
\usepackage{booktabs}
\usepackage{tcolorbox}
\tcbuselibrary{skins, breakable}
\usepackage[colorlinks=true, linkcolor=emerald, urlcolor=emerald, bookmarks=true]{hyperref}

\definecolor{emerald}{RGB}{0,102,51}
\definecolor{darkred}{RGB}{139,0,0}
\definecolor{lightgray}{RGB}{245,245,245}

\usepackage[nil]{babel}
\babelprovide[import, main]{bengali}
\babelprovide[import]{arabic}
\babelfont{rm}[Renderer=HarfBuzz,Scale=1.0]{Noto Serif Bengali}
\babelfont{sf}[Renderer=HarfBuzz]{Noto Sans Bengali}
\babelfont{tt}[Renderer=HarfBuzz]{Noto Sans Bengali}
\babelfont[arabic]{rm}[Renderer=HarfBuzz,Scale=1.1]{Noto Naskh Arabic}

% Section styling
\titleformat{\section}{\Large\bfseries\color{emerald}}{\thesection.}{0.5em}{}
\titleformat{\subsection}{\large\bfseries\color{emerald!85!black}}{\thesubsection.}{0.5em}{}
\titleformat{\subsubsection}{\normalsize\bfseries\color{emerald!70!black}}{\thesubsubsection.}{0.5em}{}
\titlespacing*{\section}{0pt}{18pt}{8pt}

% Fancy Header/Footer setup
\pagestyle{fancy}
\fancyhf{} % clear all header and footer fields
\fancyhead[L]{\color{emerald}\small\textbf{\nouppercase{\leftmark}}}
\fancyfoot[L]{\scriptsize\color{black!50}{\lat{AI}}-সংকলিত, আলেম-যাচাইকৃত নয়}
\fancyfoot[C]{\color{emerald!70!black}\thepage}
\renewcommand{\headrulewidth}{0.4pt}
\renewcommand{\headrule}{\hbox to\headwidth{\color{emerald}\leaders\hrule height \headrulewidth\hfill}}
\renewcommand{\footrulewidth}{0pt}

% Custom boxes
% 1. Ayat/Hadith Box
\newtcolorbox{ayatbox}[1][]{
    enhanced, breakable, 
    colback=emerald!5, 
    colframe=emerald, 
    boxrule=0.5pt, 
    arc=3pt, 
    left=10pt, right=10pt, top=10pt, bottom=10pt,
    #1
}

% 2. Quote/Translation Box
\newtcolorbox{quotebox}[1][]{
    enhanced, breakable,
    frame hidden,
    colback=lightgray,
    borderline west={3pt}{0pt}{emerald},
    left=15pt, right=10pt, top=10pt, bottom=10pt,
    sharp corners,
    #1
}

% 3. AI-generated notice box
\newtcolorbox{warnbox}[1][]{
    enhanced, breakable,
    colback=red!4,
    colframe=darkred,
    boxrule=0.8pt,
    arc=3pt,
    left=10pt, right=10pt, top=10pt, bottom=10pt,
    #1
}

% Commands
\newcommand{\arab}[1]{\foreignlanguage{arabic}{#1}}
\newcommand{\kib}[1]{\textcolor{emerald}{\textbf{#1}}}

% Latin (English) fallback: Noto Serif Bengali has NO Latin glyphs
\newfontfamily\latfont[Renderer=HarfBuzz]{Noto Serif}
\newcommand{\lat}[1]{{\latfont #1}}

\setlength{\parskip}{5pt}
\setlength{\parindent}{0pt}

% Fix for missing bullet in Noto Serif Bengali
\renewcommand{\labelitemi}{$\bullet$}



\begin{document}
\begin{center}{\Large\bfseries\color{emerald} আল্লাহর রহমত ও তাওবা — \lat{05}-পুড়ে-ছাই-হওয়ার-ভয়-ও-আল্লাহর-ক্ষমা}\end{center}
\vspace{1em}
\section*{ঘটনা ৫ — পুড়ে-ছাই হওয়ার ভয়: এক নেক কাজহীন ব্যক্তির অন্তরের ভয় ও ক্ষমা}

\begin{warnbox}
 \textbf{গুরুত্বপূর্ণ নোটিশ (\lat{Important} \lat{Notice}):} এই কন্টেন্টটি \lat{AI} (কৃত্রিম বুদ্ধিমত্তা) দিয়ে প্রস্তুত ও সংকলিত। \lat{AI} কঠোর সূত্র-মান নীতি মেনে শুধু নির্ভরযোগ্য উৎস থেকে তথ্য সংগ্রহ করেছে — কেবল সহীহ/হাসান হাদিস ও সহীহ সনদের আসার রাখা হয়েছে, দুর্বল সনদ বাদ দেওয়া হয়েছে। তবে এটি কোনো আলেম বা মুহাদ্দিস কর্তৃক যাচাইকৃত নয়।
\end{warnbox}

\begin{quote}
\textbf{সম্পর্কিত ফাইল:} কুরআন ফাইল — ৪০:৩ (ক্ষমা ও কঠোর শাস্তির ভারসাম্য); ব্যবহারিক গাইড — ভয় ও আশার ভারসাম্য।
\end{quote}

\medskip\hrule\medskip

\subsection*{১. সূত্র কার্ড (\lat{Source} \lat{Card})}

\begin{table}[h]\centering\small
\begin{tabular}{ll}
বিষয় & বিবরণ \\
\hline
\textbf{ঘটনা} & বনি ইসরাঈলের এক ব্যক্তি — কোনো নেক কাজ করেনি — মৃত্যুর সময় পুড়ে-ছাই হওয়ার নির্দেশ; আল্লাহ প্রশ্ন করলেন — কেন? জবাব: আপনার ভয়ে; আল্লাহ ক্ষমা করলেন \\
\textbf{বর্ণনাকারী} & আবু হুরাইরা (রাঃ) \\
\textbf{গ্রন্থ ও নম্বর} & সহীহ মুসলিম ২৭৫৬; সহীহ বুখারি ৩৪৭৮ \\
\textbf{অধ্যায়} & কিতাবুত তাওবাহ — \lat{“}রহমতের ব্যাপকতা ও তা গজবের উপর বিজয়ী\lat{”} \\
\textbf{সনদের মান} & \textbf{সহীহ} (বুখারি-মুসলিম) \\
\textbf{স্থান ও সময়} & বনি ইসরাঈলের যুগের ঘটনা — নবীজি (সা.) বর্ণনা করেছেন \\
\hline
\end{tabular}
\end{table}

\medskip\hrule\medskip

\subsection*{২. সংক্ষিপ্ত পটভূমি (\lat{Background})}

এই হাদিসটি রহমতের ব্যাপকতার অধ্যায়ে এসেছে — সেটাই এর মূল প্রসঙ্গ। ব্যক্তিটি কোনো নেক কাজ করেনি — এমনকি তাওবাও করেনি বলে বোঝা যায়; কিন্তু তার অন্তরে আল্লাহর ভয় ছিল — পরকালের হিসাবের ভয়। সেই ভয়ই শেষ পর্যন্ত তার মুক্তির কারণ হলো। হাদিসটি শেখায় — \textbf{হতাশার বিপরীত ঔষধ ভয় নয়, বরং ভয় নিজেই রহমতের দরজা খুলে দিতে পারে} — কারণ ভয়ের অর্থ অন্তর আল্লাহকে গুরুত্ব দিচ্ছে।

\medskip\hrule\medskip

\subsection*{৩. পূর্ণ ঘটনার বিশুদ্ধ বর্ণনা (\lat{The} \lat{Full} \lat{Narration})}

\begin{quote}
আবু হুরাইরা (রাঃ) থেকে — নবীজি (সা.) বলেছেন: \textbf{\lat{“}তোমাদের আগে এক ব্যক্তি ছিল, যে কোনো নেক কাজ করেনি। মৃত্যুর সময় সে তার পরিবারকে বলল: আমি মারা গেলে আমাকে পুড়িয়ে ফেলো, এরপর আমার ছাইয়ের অর্ধেক স্থলে ছড়িয়ে দাও, আর অর্ধেক সাগরে। আল্লাহর কসম! যদি আল্লাহ আমাকে (তার ক্ষমতায়) ধরতে পারেন, তবে তিনি আমাকে এমন শাস্তি দেবেন, যা তিনি দুনিয়ার কাউকে দেননি। সে মারা গেলে তার নির্দেশ মতো করা হলো। আল্লাহ স্থলকে আদেশ দিলেন — তাতে থাকা (ছাই) জড়ো করল; সাগরকে আদেশ দিলেন — তাতে থাকা (ছাই) জড়ো করল। তারপর আল্লাহ তাকে জিজ্ঞেস করলেন: কেন এমন করলে? সে বলল: হে আমার রব! আপনার ভয়েই (করেছি) — আর আপনি তো সবচেয়ে ভালো জানেন। আল্লাহ তাকে ক্ষমা করে দিলেন।\lat{”}}
\end{quote}

\medskip\hrule\medskip

\subsection*{৪. শব্দের অর্থ (\lat{Key} \lat{Words})}

\begin{table}[h]\centering\small
\begin{tabular}{ll}
শব্দ & অর্থ \\
\hline
\textbf{হাররিকূহু (\arab{حرّقوه})} & পুড়িয়ে ফেলো \\
\textbf{ইযরূ (\arab{اذروا})} & ছড়িয়ে দাও \\
\textbf{ইন কাদারা আলাইহি} & যদি আল্লাহ তার উপর ক্ষমতা পেতেন — তার ভয়ের উক্তি \\
\textbf{মিন খাশইয়াতিকা (\arab{من} \arab{خشيتك})} & আপনার ভয়ের কারণে — অন্তরের ভয় \\
\textbf{ফাগাফারা লাহু} & আল্লাহ তাকে ক্ষমা করে দিলেন \\
\hline
\end{tabular}
\end{table}

\medskip\hrule\medskip

\subsection*{৫. সংশ্লিষ্ট দলিল ও রেফারেন্স}

\begin{quote}
\textbf{\lat{“}তিনি গুনাহ ক্ষমাকারী, তাওবা গ্রহণকারী, কঠোর শাস্তিদাতা…\lat{”}} — সূরা আল-মু'মিন ৪০:৩
\end{quote}

\begin{quote}
\textbf{\lat{“}আমার রহমত আমার গজবের উপর প্রাধান্য পায়।\lat{”}} — সহীহ মুসলিম ২৭৫১
\end{quote}

\begin{quote}
\textbf{\lat{“}তোমার গুনাহ যদি আকাশের মেঘ পর্যন্ত পৌঁছে যায়, তারপর তুমি ক্ষমা চাও — আমি ক্ষমা করব।\lat{”}} — তিরমিযী ৩৫৪০ (হাসান)
\end{quote}

\subsubsection*{মূল শিক্ষা (দলিলভিত্তিক)}

\begin{enumerate}
\item \textbf{ভয়ই উদ্ধার করল:} এক নেক কাজহীন ব্যক্তি — কোনো তাওবার ভাষাও সে করেনি — শুধু অন্তরের ভয়ের কারণে ক্ষমা পেল; আল্লাহ অন্তর দেখেন — পাপীর অন্তরে আল্লাহর ভয় থাকলেই আশা আছে।
\item \textbf{ভয় ও আশার মিলন:} সে ভেবেছিল শাস্তি পাবেই (ভয়) — কিন্তু আল্লাহ দেখলেন সেই ভয়ই তার অন্তরে আল্লাহর প্রতি গুরুত্বের প্রমাণ (আশা); রহমত গজবের উপর বিজয়ী (মুসলিম ২৭৫১)।
\item \textbf{হতাশ নয় — ভীত:} হাদিসের শিক্ষা — হতাশা নয়; ভয়কে আল্লাহর দিকে ফেরার সোপান বানানো; ভয় যত গভীর, রহমতের দরজা তত কাছে।
\item \textbf{সতর্কতা:} এর মানে পাপে থাকা ভালো — তা নয়; বরং যারা অতীত পাপে ভেঙে পড়েছে, তাদের বলা — অন্তরের ভয়টুকু থাকলেই আল্লাহর রহমতের দরজা বন্ধ নয়।
\end{enumerate}
\medskip\hrule\medskip

\subsection*{৬. রেফারেন্স}

\begin{itemize}
\item সহীহ মুসলিম ২৭৫৬ (কিতাবুত তাওবাহ)
\item সহীহ বুখারি ৩৪৭৮ (কিতাবুল আম্বিয়া)
\item কুরআন: মু'মিন ৪০:৩
\item ব্যাখ্যা: শারহু নববী — রহমতের ব্যাপকতার অধ্যায়
\end{itemize}

\end{document}
